\documentclass[11pt]{article}
\usepackage[margin=1in]{geometry}
\usepackage{amsmath,amssymb,bm,microtype}
\usepackage[colorlinks=true,allcolors=blue]{hyperref}
\newcommand{\Kern}{\mathcal{K}}
\newcommand{\beq}{\mathrm{eq}}
\newcommand{\bint}{\mathrm{int}}
\newcommand{\eff}{\mathrm{eff}}
\title{Detailed Balance for Effective One-Dimensional Reaction Kernels}
\author{}
\date{}
\begin{document}
\maketitle

\section{Reviewer comment and proposed response}
\subsection*{Reviewer comment}
\begin{quote}
A related question concerns the use of detailed-balance emission rates
together with effective sink-strength reaction kernels. While Eq.~(70) is
standard for conventional second-order absorption reactions, the manuscript
also relies on reaction kernels derived from 1D-mobile defect/fixed-sink
formulations. It is unclear whether the same detailed-balance arguments remain
applicable in such third-order reactions. At the very least, the assumptions
and limitations underlying this approximation should be discussed.
\end{quote}

\subsection*{Proposed response}
\begin{quote}
We thank the reviewer for raising this important point. We agree that the
detailed-balance relation conventionally used for a binary absorption reaction
cannot be transferred without qualification to an effective,
concentration-dependent sink-strength kernel.

The pure-1D absorption law is not interpreted here as an elementary three-body
reaction. It results from averaging the first-passage motion of a 1D-mobile
cluster over the entire fixed-sink population. Its apparent third-order
dependence arises because one factor of sink concentration specifies the
family responsible for capture, whereas the total macroscopic cross section
determines the mean distance between intersections of the glide trajectory and
all sinks \cite{Borodin1998,Barashev2001,Trinkaus2002}.

Detailed balance applies rigorously to the underlying reversible elementary
transition \(A+B\rightleftharpoons C\), rather than automatically to its
coarse-grained multisink closure. At equilibrium,
\[
K^{\eff}_{AB}(\bm c^{\beq})c_A^{\beq}c_B^{\beq}
=\alpha_C^{\eff}(\bm c^{\beq})c_C^{\beq}.
\]
Thus, if the concentration-dependent 1D absorption kernel is retained in the
equilibrium relation, the thermodynamically conjugate effective emission
coefficient must contain the same environmental dependence. Combining an
instantaneous 1D sink-strength kernel with a concentration-independent
emission rate obtained from the conventional binary formula is not an exact
detailed-balance construction.

In the present model, the detailed-balance emission expression is therefore
regarded as a dilute-medium, local-equilibrium approximation. Emission is
assumed to be governed by the binding free energy of the emitting cluster,
whereas the multisink transport correction affects capture without changing
the intrinsic detachment event. The revised manuscript states this assumption
and its limitations. The approximation may lose thermodynamic consistency
when sink populations change substantially, when 1D transport produces strong
spatial correlations, or when the system approaches thermal equilibrium.
\end{quote}

If emission is demonstrably small in the simulated regime, one may add:
\begin{quote}
Because the calculations address a strongly driven irradiation regime in
which production and absorption fluxes are much greater than the relevant
thermally activated emission fluxes, this approximation does not control the
predicted microstructure. Nevertheless, its thermodynamic status and range of
validity are now stated explicitly.
\end{quote}

\section{Detailed balance for a general reaction}
Consider
\begin{equation}
 \sum_i\nu_{ir}^{-}X_i\rightleftharpoons
 \sum_i\nu_{ir}^{+}X_i,\qquad
 \Delta\nu_{ir}=\nu_{ir}^{+}-\nu_{ir}^{-}.
\end{equation}
For ideal mass-action kinetics,
\begin{equation}
 J_r^{+}=k_r^{+}\prod_i a_i^{\nu_{ir}^{-}},\qquad
 J_r^{-}=k_r^{-}\prod_i a_i^{\nu_{ir}^{+}},
\end{equation}
where \(a_i\) is a dimensionless activity. Microscopic reversibility requires
\(J_r^{+,\beq}=J_r^{-,\beq}\) for every reversible channel. Consequently,
\begin{equation}
 \boxed{\frac{k_r^{+}}{k_r^{-}}
 =\prod_i(a_i^{\beq})^{\Delta\nu_{ir}}\equiv K_r^{\beq}},
 \qquad
 \boxed{k_r^{-}=\frac{k_r^{+}}{K_r^{\beq}}}.
 \label{eq:general-db}
\end{equation}
When dimensional concentrations replace activities, powers of the standard
concentration must be retained because forward and reverse constants generally
have different dimensions \cite{Hill1989,Qian2022}.

\section{Conventional binary absorption and emission}
For \(X_n+X_1\rightleftharpoons X_{n+1}\), write
\[
 J_n^{+}=\beta_nc_nc_1,\qquad
 J_{n+1}^{-}=\alpha_{n+1}c_{n+1}.
\]
Detailed balance gives
\begin{equation}
 \boxed{\alpha_{n+1}
 =\beta_n\frac{c_n^{\beq}c_1^{\beq}}{c_{n+1}^{\beq}}}.
 \label{eq:binary-db}
\end{equation}
If
\[
 \frac{c_{n+1}^{\beq}}{c_n^{\beq}c_1^{\beq}}
 =g_{n+1}\exp\left(\frac{\Delta G_{n+1}^{b}}{k_{\rm B}T}\right),
\]
then
\begin{equation}
 \boxed{\alpha_{n+1}
 =\frac{\beta_n}{g_{n+1}}
 \exp\left(-\frac{\Delta G_{n+1}^{b}}{k_{\rm B}T}\right)}.
\end{equation}
Here \(g_{n+1}\) contains standard-state and configurational-degeneracy
factors. Replacing binding free energy by binding energy assumes that the
relevant entropy is negligible or included in the prefactor.

\section{Pure-1D absorption as an effective higher-order law}
Let 1D-mobile species \(A\) be captured by fixed clusters of family \(B\).
Define
\begin{equation}
 \Sigma_A=\sum_j\sigma_{Aj}N_j.
\end{equation}
The family-resolved sink strength and per-\(A\)-cluster frequency are
\begin{equation}
 k_{(1),A\to B}^{2}=6\sigma_{AB}N_B\Sigma_A,\qquad
 \Kern^{(1)}_{A\to B}
 =6D_A\sigma_{AB}N_B\sum_j\sigma_{Aj}N_j.
\end{equation}
The rate density is
\begin{equation}
 \boxed{J_{A+B}^{+}
 =6D_A\sigma_{AB}N_AN_B\sum_j\sigma_{Aj}N_j}
 =\sum_j6D_A\sigma_{AB}\sigma_{Aj}N_AN_BN_j.
 \label{eq:1d-rate}
\end{equation}
This resembles \(A+B+X_j\to C+X_j\), but \(X_j\) is not normally a third
defect in the local capture event. Its population determines the ensemble
spacing between intersections of the 1D path and sinks. The third factor is
environmental and geometrical, generated by coarse graining
\cite{Borodin1998,Barashev2001,Adjanor2022}.

For one sink family,
\begin{equation}
 \boxed{J_{A+B}^{+}=6D_A\sigma_{AB}^{2}N_AN_B^2}.
\end{equation}
With \(\sigma_{AB}=\pi R_{AB}^{2}\),
\begin{equation}
 J_{A+B}^{+}=6\pi^2D_AR_{AB}^{4}N_AN_B^2,
\end{equation}
displaying the effective third-order and \(R^4\) dependences.

\section{Detailed balance for the effective 1D closure}
For \(A+B\rightleftharpoons C\), write
\begin{equation}
 J^{+}=K_{AB}^{\eff}(\bm N)N_AN_B,\qquad
 K_{AB}^{\eff}(\bm N)=6D_A\sigma_{AB}\Sigma_A,
\end{equation}
and \(J^{-}=\alpha_C^{\eff}(\bm N)N_C\). At equilibrium,
\begin{equation}
 \boxed{\alpha_C^{\eff}(\bm N^{\beq})
 =K_{AB}^{\eff}(\bm N^{\beq})
 \frac{N_A^{\beq}N_B^{\beq}}{N_C^{\beq}}}.
 \label{eq:effective-db}
\end{equation}
Defining
\[
 K_{\rm th}=\frac{N_C^{\beq}}{N_A^{\beq}N_B^{\beq}},
\]
gives
\begin{equation}
 \boxed{\alpha_C^{\eff}(\bm N^{\beq})
 =\frac{K_{AB}^{\eff}(\bm N^{\beq})}{K_{\rm th}}}.
\end{equation}
This guarantees zero net flux at that equilibrium; it does not define an
intrinsic detachment frequency away from that state.

\subsection{Instantaneously consistent coarse-grained closure}
One can impose
\[
 \alpha_C^{\eff}(\bm N)=K_{AB}^{\eff}(\bm N)/K_{\rm th}
\]
at every instant. Then
\begin{equation}
 \boxed{J_{\rm net}=K_{AB}^{\eff}(\bm N)
 \left(N_AN_B-\frac{N_C}{K_{\rm th}}\right)}.
 \label{eq:instantaneous-net}
\end{equation}
This preserves the equilibrium manifold, but makes the effective emission
frequency depend on the external sink population through \(\Sigma_A\). It is
a consistent coarse-grained reverse flux, not the intrinsic frequency at
which \(C\) emits \(A\).

\section{Formulation for genuine higher-order reactions}
For
\[
 \sum_i\alpha_iX_i\rightleftharpoons\sum_i\beta_iX_i,
\]
the generalized mass-action flux is
\begin{equation}
 J_{\rm net}
 =k^+\prod_i a_i^{\alpha_i}
 -k^-\prod_i a_i^{\beta_i},\qquad
 \frac{k^+}{k^-}
 =\prod_i(a_i^{\beq})^{\beta_i-\alpha_i}.
\end{equation}
Detailed balance therefore applies to any molecularity, provided the forward
and reverse laws represent the same elementary channel and use consistent
activities.

For the formal catalytic reaction
\[
 A+B+S\rightleftharpoons C+S,
\]
\[
 J^{+}=k^{+}a_Aa_Ba_S,\qquad J^{-}=k^{-}a_Ca_S.
\]
The catalyst cancels from the equilibrium constant,
\begin{equation}
 \boxed{\frac{k^{+}}{k^{-}}
 =\frac{a_C^{\beq}}{a_A^{\beq}a_B^{\beq}}},
\end{equation}
but multiplies both nonequilibrium fluxes:
\[
 J_{\rm net}=a_S(k^{+}a_Aa_B-k^{-}a_C).
\]
Thus a formally reversible multisink closure is
\begin{equation}
 \boxed{J_{\rm net}
 =6D_A\sigma_{AB}\Sigma_A
 \left(N_AN_B-\frac{N_C}{K_{\rm th}}\right)}.
\end{equation}
This is a reversible coarse-grained equation, but its reverse term is not an
intrinsic evaporation rate.

\section{Recommended physical formulation}
The cleanest treatment separates intrinsic detachment from transport and
recapture:
\begin{equation}
 \boxed{\alpha_C^{\bint}
 =\nu_C\exp\left(-\frac{\Delta G_C^\ddagger}{k_{\rm B}T}\right)}.
 \label{eq:intrinsic}
\end{equation}
If
\[
 \Delta G_C^\ddagger\simeq\Delta G_A^m+\Delta G_C^b,
\]
then
\begin{equation}
 \boxed{\alpha_C^{\bint}\simeq\nu_C
 \exp\left[-\frac{\Delta G_A^m+\Delta G_C^b}{k_{\rm B}T}\right]}.
\end{equation}
The intrinsic event should not depend on unrelated external sinks. After
emission, the released defect undergoes 1D, 1DR, or 3D transport and possible
recapture, described separately by the appropriate transport kernel.

\section{Modeling alternatives and limitations}
\begin{enumerate}
 \item \textbf{Mechanistic formulation (recommended):} compute intrinsic
 emission from the detachment barrier and use the 1D or 1DR kernel only for
 subsequent capture and recapture.
 \item \textbf{Reversible coarse-grained formulation:} multiply the reverse
 flux by the same concentration-dependent transport factor as the forward
 flux, as in Eq.~\eqref{eq:instantaneous-net}. This preserves equilibrium but
 makes the reverse coefficient an effective environmental rate.
 \item \textbf{Dilute nonequilibrium approximation:} retain conventional
 binary detailed-balance emission while using the effective 1D absorption
 kernel. This may be adequate when emission is demonstrably negligible, but
 it is not an exact detailed-balance construction.
\end{enumerate}
The third approximation may fail when sink populations evolve strongly,
spatial correlations invalidate the homogeneous closure, rotation or trapping
changes migration dimensionality, emission and recapture occur on comparable
time scales, or the system approaches thermal equilibrium.

\section{Suggested manuscript insertion}
\begin{quote}
\textbf{Detailed balance for concentration-dependent kernels.}
The detailed-balance relation conventionally used for
\(X_n+X_1\rightleftharpoons X_{n+1}\) is exact when the forward and reverse
terms represent the same pair of elementary transitions. The sink-strength
expression for pure 1D migration has a different status:
\[
K^{\eff}_{n,1}(\bm N)=6D_1\sigma_{1n}\sum_j\sigma_{1j}N_j,\qquad
J^+_{n,1}=K^{\eff}_{n,1}(\bm N)N_nN_1.
\]
Its apparent higher-order dependence is not an elementary three-defect
collision. The additional factor describes how the sink population determines
the mean distance between intersections of a 1D trajectory and absorbing
sinks \cite{Borodin1998,Barashev2001,Trinkaus2002}.

At equilibrium, a coarse-grained reverse coefficient must satisfy
\[
\alpha^{\eff}_{n+1}(\bm N^{\beq})
=K^{\eff}_{n,1}(\bm N^{\beq})
\frac{N_n^{\beq}N_1^{\beq}}{N_{n+1}^{\beq}}.
\]
Thus, applying the conventional binary detailed-balance expression with an
instantaneous concentration-dependent 1D kernel is not an exact microscopic
construction unless the reverse term contains the same environmental
dependence. Here thermal emission is treated as an intrinsic unimolecular
detachment process governed by cluster binding free energy, whereas the 1D
sink-strength factor describes subsequent transport-controlled capture. Use
of the binary expression with an effective 1D kernel is therefore a
dilute-medium, local-equilibrium approximation. It may become inaccurate when
sink populations evolve strongly, spatial correlations are important, or the
system approaches thermal equilibrium.
\end{quote}

\section*{References}
\begin{thebibliography}{99}
\bibitem{Smoluchowski1917}
M. von Smoluchowski, ``Versuch einer mathematischen Theorie der
Koagulationskinetik kolloider L{\"o}sungen,''
\emph{Z. Phys. Chem.} \textbf{92} (1917) 129--168.

\bibitem{Waite1957}
T. R. Waite, ``Theoretical treatment of the kinetics of diffusion-limited
reactions,'' \emph{Phys. Rev.} \textbf{107} (1957) 463--470,
\href{https://doi.org/10.1103/PhysRev.107.463}{doi:10.1103/PhysRev.107.463}.

\bibitem{GoseleSeeger1976}
U. G{\"o}sele and A. Seeger, ``Theory of bimolecular reaction rates limited
by anisotropic diffusion,'' \emph{Philos. Mag.} \textbf{34} (1976) 177--193,
\href{https://doi.org/10.1080/14786437608223764}
{doi:10.1080/14786437608223764}.

\bibitem{Borodin1998}
V. A. Borodin, ``Rate theory for one-dimensional diffusion,''
\emph{Physica A} \textbf{260} (1998) 467--478,
\href{https://doi.org/10.1016/S0378-4371(98)00338-0}
{doi:10.1016/S0378-4371(98)00338-0}.

\bibitem{Barashev2001}
A. V. Barashev, S. I. Golubov, and H. Trinkaus, ``Reaction kinetics of
glissile interstitial clusters in a crystal containing voids and
dislocations,'' \emph{Philos. Mag. A} \textbf{81} (2001) 2515--2532,
\href{https://doi.org/10.1080/01418610108217159}
{doi:10.1080/01418610108217159}.

\bibitem{Trinkaus2002}
H. Trinkaus, H. L. Heinisch, A. V. Barashev, S. I. Golubov, and B. N. Singh,
``1D to 3D diffusion-reaction kinetics of defects in crystals,''
\emph{Phys. Rev. B} \textbf{66} (2002) 060105(R),
\href{https://doi.org/10.1103/PhysRevB.66.060105}
{doi:10.1103/PhysRevB.66.060105}.

\bibitem{HuangMarian2019}
S. Huang and J. Marian, ``Rates of diffusion-controlled reactions for
one-dimensionally moving species in 3D space,''
\emph{Philos. Mag.} \textbf{99} (2019) 2562--2583,
\href{https://doi.org/10.1080/14786435.2019.1632930}
{doi:10.1080/14786435.2019.1632930}.

\bibitem{Adjanor2022}
G. Adjanor, ``Complete characterization of sink strengths for mutually
1D-mobile defect clusters: Extension to diffusion-anisotropy analog cases,''
\emph{J. Nucl. Mater.} \textbf{570} (2022) 153970,
\href{https://doi.org/10.1016/j.jnucmat.2022.153970}
{doi:10.1016/j.jnucmat.2022.153970}.

\bibitem{Hill1989}
T. L. Hill, \emph{Free Energy Transduction and Biochemical Cycle Kinetics},
Springer-Verlag, New York, 1989.

\bibitem{Qian2022}
H. Qian, ``Statistical chemical thermodynamics and energetic behavior of
molecular systems: From equilibrium to nonequilibrium,''
\emph{J. Phys. Chem. B} \textbf{126} (2022).
\end{thebibliography}
\end{document}
